\chapter{Slownik pojec}

\section*{Argument}
Wartosc przekazywana do funkcji podczas jej wywolania.

\section*{Chunk}
Podstawowy blok organizacyjny MoonChunk. Moze zawierac strony, metadata, funkcje,
petle, warunki i inne instrukcje.

\section*{Content}
Blok zawierajacy HTML, ktory trafi do glownej tresci strony.

\section*{Deklaracja}
Moment utworzenia nowej nazwy, np. przez \mc{let}, \mc{const} albo \mc{global}.

\section*{Diagnostyka}
Komunikaty o bledach i problemach pojawiajacych sie podczas parsowania lub wykonania.

\section*{Global}
Wspolna wartosc deklarowana w bloku \mc{env}, dostepna w szerszym zakresie niz zwykle
zmienne lokalne.

\section*{HTML}
Jezyk znacznikow opisujacy strukture dokumentu strony internetowej.

\section*{Import}
Mechanizm pobierania chunkow z innego pliku \mc{.mncnk}.

\section*{Include}
Mechanizm dolaczania jednego chunku do drugiego przez \mc{@include}.

\section*{Interpolacja}
Wstawianie dynamicznej wartosci do tekstu lub HTML-a, np. przez zapis \mc{\{title\}}.

\section*{JSON}
Popularny format zapisu danych, z ktorego MoonChunk moze wczytywac informacje przez
\mc{data(...)}.

\section*{Metadata}
Dodatkowe informacje o dokumencie HTML, zwykle trafiajace do sekcji \mc{<head>} albo
sterujace ogolna struktura strony.

\section*{Page}
Instrukcja opisujaca pojedyncza strone do wygenerowania.

\section*{Parametr}
Nazwa zdefiniowana w deklaracji funkcji, do ktorej trafia argument przy wywolaniu.

\section*{Rekurencja}
Sytuacja, w ktorej funkcja wywoluje sama siebie.

\section*{Rzutowanie}
Swiadome przeksztalcenie lub interpretacja wartosci jako innego typu.

\section*{Scope / zasieg}
Obszar programu, w ktorym dana nazwa jest widoczna i dostepna.

\section*{Static site}
Serwis zbudowany z gotowych plikow HTML, CSS i JavaScript, bez potrzeby dynamicznego
renderowania po stronie serwera przy kazdym wejsciu uzytkownika.

\section*{Typ}
Informacja o rodzaju wartosci, np. liczba, tekst, wartosc logiczna, tablica czy obiekt.

\section*{Warunek}
Wyrazenie logiczne, ktore decyduje, czy dany blok kodu ma zostac wykonany.

\section*{Wyrazenie}
Fragment kodu, ktory po obliczeniu daje wartosc.
